% Geraden – bank/geraden/ (zusammenbau v0.4, Heft)
\einheitenkopf[t3]{Geraden}
\setcounter{aufgabe}{25}

% Anlauf (ohne Prüfkennung)
\begin{aufgabe}{Geradengleichung aufstellen und lesen – Anlauf}
\begin{teile}
\teil Gegeben ist $g\colon \vec x = (4 | -1 | 2) + t \cdot (3 | 0 | -2)$. Welcher Vektor gibt die Richtung von $g$ an? \kreuz{$(4 | -1 | 2)$} \\ \kreuz{$(3 | 0 | -2)$} \\ \kreuz{$(7 | -1 | 0)$}
\teil Stelle eine Gleichung der Geraden $g$ durch $A(2 | 1 | 3)$ und $B(4 | 4 | 2)$ auf.
\end{teile}
\end{aufgabe}

\begin{aufgabe}{Geradengleichung aufstellen und lesen – Prüfungsaufgaben}
\begin{teile}
\teil Ein Skilift führt geradlinig von der Talstation $T(2 | -4 | 1)$ zur Bergstation $B(-6 | 20 | 13)$; 1 LE entspricht 10 m. Gib eine Gleichung der Liftstrecke an und beschreibe, was sie im Sachzusammenhang darstellt. \hfill (Abitur 2023 GK)
\teil Eine Seilrutsche wird durch $\vec x = (0 | 5 | 12) + \lambda \cdot (15 | -10 | -10)$ mit $\lambda \in [0; 1]$ beschrieben; 1 LE entspricht 1 m, die $x_1x_2$-Ebene ist der Boden. Beschreibe die Bedeutung der Gleichung im Sachzusammenhang. \hfill (Abitur 2023 GK)
\teil Ein Förderband wird durch $\vec x = (3 | 1 | 0) + \mu \cdot (6 | 3 | 6)$ mit $0 \le \mu \le 1$ beschrieben; 1 LE entspricht 1 m, die $x_1x_2$-Ebene ist der Boden. Beschreibe die Bedeutung der Gleichung im Sachzusammenhang. \hfill (Abitur 2023 GK)
\teil Gegeben ist $g\colon \vec x = (3 | 1 | 2) + r \cdot (1 | 3 | -2)$, $r \in \mathbb{R}$. Gib eine Gleichung einer Geraden $p$ an, die echt parallel zu $g$ verläuft, und eine Gleichung einer Geraden $q$, die $g$ senkrecht schneidet. \hfill (Abitur 2019 GK)
\teil Gegeben ist $g\colon \vec x = (-2 | 4 | 1) + t \cdot (2 | -1 | 2)$, $t \in \mathbb{R}$. Gib eine Gleichung einer Geraden $p$ an, die echt parallel zu $g$ verläuft, und eine Gleichung einer Geraden $q$, die $g$ senkrecht schneidet. \hfill (Abitur 2019 GK)
\teil Gegeben ist $g\colon \vec x = (1 | -3 | 5) + r \cdot (3 | 1 | -1)$, $r \in \mathbb{R}$. Gib eine Gleichung einer Geraden $p$ an, die echt parallel zu $g$ verläuft, und eine Gleichung einer Geraden $q$, die $g$ senkrecht schneidet. \hfill (Abitur 2019 GK)
\teil Die Geraden $g$ und $h$ verlaufen parallel mit dem Richtungsvektor $(1 | 4 | 2)$; $P(2 | -1 | 3)$ liegt auf $g$, $Q(10 | 3 | -1)$ auf $h$. Der Punkt $T$ liegt auf der Strecke $PQ$ mit $|TQ| = 3 \cdot |PT|$. Gib eine Gleichung der Geraden an, die durch $T$ parallel zu $g$ und $h$ verläuft. \hfill (Abitur 2022 GK)
\teil Die Geraden $g$ und $h$ verlaufen parallel mit dem Richtungsvektor $(2 | -1 | 1)$; $P(-3 | 0 | 5)$ liegt auf $g$, $Q(3 | 6 | -1)$ auf $h$. Der Punkt $T$ liegt auf der Strecke $PQ$ mit $|TQ| = 2 \cdot |PT|$. Gib eine Gleichung der Geraden an, die durch $T$ parallel zu $g$ und $h$ verläuft. \hfill (Abitur 2022 GK)
\teil Die Geraden $g$ und $h$ verlaufen parallel mit dem Richtungsvektor $(1 | 1 | -2)$; $P(0 | 4 | 1)$ liegt auf $g$, $Q(8 | 0 | 5)$ auf $h$. Der Punkt $T$ liegt auf der Strecke $PQ$ mit $|PT| = 3 \cdot |TQ|$. Gib eine Gleichung der Geraden an, die durch $T$ parallel zu $g$ und $h$ verläuft. \hfill (Abitur 2022 GK)
\teil Der Punkt $P(3 | 2 | 2)$ liegt in der Ebene $E\colon \vec x = (2 | -1 | 1) + r \cdot (1 | 0 | 2) + t \cdot (0 | 3 | -1)$. Die Gerade $g$ liegt in $E$, geht durch $P$ und hat keinen Schnittpunkt mit der $x_1x_2$-Ebene. Gib eine Gleichung von $g$ an. \hfill (Abitur 2021 GK)
\teil Der Punkt $P(4 | 1 | 6)$ liegt in der Ebene $E\colon \vec x = (0 | 2 | 3) + r \cdot (2 | 1 | -1) + t \cdot (1 | -1 | 2)$. Die Gerade $g$ liegt in $E$, geht durch $P$ und hat keinen Schnittpunkt mit der $x_1x_2$-Ebene. Gib eine Gleichung von $g$ an. \hfill (Abitur 2021 GK)
\end{teile}
\end{aufgabe}

% Anlauf (ohne Prüfkennung)
\begin{aufgabe}{Punktprobe und Punkte auf der Geraden – Anlauf}
\begin{teile}
\teil Für $g\colon \vec x = (2 | 5 | 1) + t \cdot (0 | 3 | 4)$ und $P(2 | 11 | 9)$: Aus welcher Koordinate lässt sich $t$ bestimmen? \kreuz{nur aus der ersten} \\ \kreuz{aus der zweiten oder der dritten} \\ \kreuz{aus keiner}
\teil Prüfe, ob $P(7 | 5 | 8)$ auf $g\colon \vec x = (1 | 2 | -1) + r \cdot (2 | 1 | 3)$ liegt.
\end{teile}
\end{aufgabe}

\begin{aufgabe}{Punktprobe und Punkte auf der Geraden – Prüfungsaufgaben}
\begin{teile}
\teil Gegeben sind $g\colon \vec x = (5 | -1 | 2) + t \cdot (3 | 0 | -2)$, $t \in \mathbb{R}$, und $A(2 | 3 | 4)$. Begründe, dass $A$ nicht auf $g$ liegt. \hfill (Abitur 2023 GK)
\teil Gegeben sind $h\colon \vec x = (0 | 4 | 6) + r \cdot (2 | -3 | 0)$, $r \in \mathbb{R}$, und $B(4 | -2 | 5)$. Begründe, dass $B$ nicht auf $h$ liegt. \hfill (Abitur 2023 GK)
\teil Gegeben sind $k\colon \vec x = (-3 | 2 | 1) + t \cdot (0 | 1 | 4)$, $t \in \mathbb{R}$, und $C(1 | 5 | 13)$. Begründe, dass $C$ nicht auf $k$ liegt. \hfill (Abitur 2023 GK)
\teil Der Punkt $B(-2 | y_B | z_B)$ liegt auf $g\colon \vec x = (1 | 3 | 2) + r \cdot (-1 | 2 | -3)$. Bestimme $y_B$ und $z_B$. \hfill (Abitur 2019 GK) \\ $y_B$ = \leerfeld , $z_B$ = \leerfeld
\teil Der Punkt $P(x_P | 2 | z_P)$ liegt auf $g\colon \vec x = (0 | -1 | 4) + t \cdot (2 | 1 | -2)$. Bestimme $x_P$ und $z_P$. \hfill (Abitur 2019 GK) \\ $x_P$ = \leerfeld , $z_P$ = \leerfeld
\teil Der Punkt $Q(x_Q | y_Q | -3)$ liegt auf $h\colon \vec x = (3 | -2 | 0) + t \cdot (-2 | 4 | 1)$. Bestimme $x_Q$ und $y_Q$. \hfill (Abitur 2019 GK) \\ $x_Q$ = \leerfeld , $y_Q$ = \leerfeld
\teil Die Gerade $g$ verläuft durch $A(3 | 1 | 2)$ mit dem Richtungsvektor $\vec u = (2 | -1 | 2)$. Bestimme die Koordinaten eines Punkts auf $g$, der von $A$ den Abstand 9 hat. \hfill (Abitur 2019 GK)
\teil Die Gerade $g$ verläuft durch $A(0 | 4 | -1)$ mit dem Richtungsvektor $\vec u = (4 | 0 | -3)$. Bestimme die Koordinaten eines Punkts auf $g$, der von $A$ den Abstand 10 hat. \hfill (Abitur 2019 GK)
\teil Die Gerade $g$ verläuft durch $A(-2 | 1 | 3)$ mit dem Richtungsvektor $\vec u = (1 | 2 | 2)$. Bestimme die Koordinaten eines Punkts auf $g$, der von $A$ den Abstand $4{,}5$ hat. \hfill (Abitur 2019 GK)
\teil Eine Hecke verläuft geradlinig von $E(8 | 0 | 0)$ nach $F(0 | 8 | 0)$ (1 LE = 1 m). Weise nach, dass der Schattenpunkt $T(2 | 6 | 0)$ auf der Strecke $EF$ liegt. \hfill (Abitur 2018 GK)
\teil Ein Zaun verläuft geradlinig von $E(6 | -2 | 0)$ nach $F(2 | 2 | 0)$ (1 LE = 1 m). Entscheide, ob der Schattenpunkt $T(-3 | 7 | 0)$ auf dem Zaun $EF$ liegt. \hfill (Abitur 2018 GK)
\teil Ein Stahlseil ist geradlinig von $A(1 | 1 | 4)$ nach $B(7 | 4 | 1)$ gespannt (1 LE = 1 m). Weise nach, dass der Punkt $L(5 | 3 | 2)$, an dem eine Lampe hängt, auf dem Seil liegt. \hfill (Abitur 2018 GK)
\end{teile}
\end{aufgabe}

\begin{aufgabe}{Punktprobe und Punkte auf der Geraden – Prüfungsaufgaben – weiter}
\begin{teile}
\teil Weise nach, dass die Punkte $A(2 | -1 | 3)$, $B(4 | 0 | 1)$ und $C(8 | 2 | -3)$ auf einer Geraden liegen. \hfill (Abitur 2020 GK)
\teil Weise nach, dass die Punkte $A(0 | 5 | 2)$, $B(1 | 3 | 3)$ und $C(-2 | 9 | 0)$ auf einer Geraden liegen. \hfill (Abitur 2020 GK)
\teil Untersuche, ob die Punkte $A(1 | 0 | 2)$, $B(3 | 1 | 5)$ und $C(7 | 3 | 12)$ auf einer Geraden liegen. \hfill (Abitur 2020 GK)
\teil Ein Sessellift führt geradlinig von $A(1 | 20 | 2)$ nach $E(7 | 2 | 11)$; die Talstation liegt im Ursprung, die $x_1x_2$-Ebene ist waagerecht, 1 LE entspricht 5 m. Wie hoch über der Talstation ist der Lift an der Stelle, die 84 m vom Anfang $A$ entfernt ist? \hfill (Abitur 2023 GK) \\ \leerfeld[m]
\teil Eine Materialseilbahn führt geradlinig von $A(0 | 0 | 4)$ nach $E(8 | 4 | 12)$; der Ursprung liegt auf Höhe der Talstation, die $x_1x_2$-Ebene ist waagerecht, 1 LE entspricht 10 m. Wie hoch über der Talstation ist die Bahn 30 m nach dem Anfang $A$? \hfill (Abitur 2023 GK) \\ \leerfeld[m]
\teil Eine Standseilbahn fährt geradlinig von $A(3 | 1 | 0)$ nach $E(-1 | 9 | 8)$; die Talstation liegt im Ursprung, die $x_1x_2$-Ebene ist waagerecht, 1 LE entspricht 20 m. Wie hoch über der Talstation ist die Bahn 180 m nach dem Anfang $A$? \hfill (Abitur 2023 GK) \\ \leerfeld[m]
\teil Ein Prisma hat die Ecken $A(5 | 0 | 0)$ und $G(1 | 4 | 6)$. Gib eine Gleichung der Geraden durch $A$ und $G$ an und weise nach, dass der Punkt $S(3 | 2 | 5)$ nicht auf ihr liegt. \hfill (Abitur 2026 GK)
\teil Ein Quader hat die Ecken $A(0 | 6 | 1)$ und $G(4 | 0 | 4)$. Gib eine Gleichung der Geraden durch $A$ und $G$ an und weise nach, dass der Punkt $S(2 | 3 | 2)$ nicht auf ihr liegt. \hfill (Abitur 2026 GK)
\teil Ein Prisma hat die Ecken $A(2 | 0 | 0)$ und $G(-2 | 6 | 6)$. Gib eine Gleichung der Geraden durch $A$ und $G$ an und untersuche, ob der Punkt $S(0 | 3 | 3)$ auf ihr liegt. \hfill (Abitur 2026 GK)
\end{teile}
\end{aufgabe}

\begin{aufgabe}{Punkt auf einer Geraden mit vorgegebenem Abstand zu einem festen Punkt über eine quadratische Gleichung bestimmen – Prüfungsaufgaben}
\begin{teile}
\teil Von einem Mast in $F(4 | -3 | 5)$ soll ein $7{,}5$ m langes Seil zu einem Punkt der Laufschiene $g\colon \vec x = (0 | 1 | 3) + t \cdot (1 | 2 | 2)$ gespannt werden (1 LE = 1 m). Bestimme alle möglichen Befestigungspunkte auf $g$. \hfill (Abitur 2017 GK)
\end{teile}
\end{aufgabe}

% Anlauf (ohne Prüfkennung)
\begin{aufgabe}{Lage als Anhang – Anlauf}
\begin{teile}
\teil Die Richtungsvektoren der Geraden $g$ und $h$ sind Vielfache voneinander. Was ist als Nächstes zu prüfen? \kreuz{ob ein Stützpunkt von $g$ auf $h$ liegt} \\ \kreuz{ob sich die Geraden schneiden} \\ \kreuz{nichts – die Geraden sind identisch}
\teil Gegeben sind $g\colon \vec x = (2 | -1 | 3) + r \cdot (1 | 3 | 0)$ und $h\colon \vec x = (2 | -1 | 3) + t \cdot (3 | 1 | 0)$. Begründe, dass $g$ und $h$ nicht identisch sind.
\teil Ein Tunnel verläuft geradlinig durch $A(0 | 30 | 5)$ und $B(30 | 50 | 8)$, ein zweiter durch $C(10 | 0 | 2)$ und $D(70 | 40 | 8)$ (1 LE = 10 m). Prüfe, ob die beiden Tunnelachsen parallel sind.
\end{teile}
\end{aufgabe}

\begin{aufgabe}{Lage als Anhang – Prüfungsaufgaben}
\begin{teile}
\teil Die Gerade $g$ durch $P(1 | 0 | 2)$ und die Gerade $h$ durch $Q(5 | -2 | 7)$ haben beide den Richtungsvektor $(2 | -1 | 3)$. Weise nach, dass $g$ und $h$ nicht identisch sind. \hfill (Abitur 2022 GK)
\teil Die Gerade $g$ durch $P(-2 | 3 | 0)$ und die Gerade $h$ durch $Q(0 | 4 | 4)$ haben beide den Richtungsvektor $(1 | 0 | 2)$. Weise nach, dass $g$ und $h$ nicht identisch sind. \hfill (Abitur 2022 GK)
\teil Die Gerade $g$ durch $P(3 | 1 | -1)$ und die Gerade $h$ durch $Q(9 | -2 | 4)$ haben beide den Richtungsvektor $(-2 | 1 | -2)$. Weise nach, dass $g$ und $h$ nicht identisch sind. \hfill (Abitur 2022 GK)
\teil Ein Würfel hat die Ecken $A(5 | 0 | 0)$, $B(5 | 5 | 0)$, $C(0 | 5 | 0)$ und $G(0 | 5 | 5)$. Gib eine Gleichung der Geraden $h$ durch $B$ und $G$ an und begründe, dass $h$ windschief zur Geraden $AC$ ist. \hfill (Abitur 2021 GK)
\teil Eine Pyramide hat die Grundfläche $A(0 | 0 | 0)$, $B(8 | 0 | 0)$, $C(8 | 8 | 0)$, $D(0 | 8 | 0)$ und die Spitze $S(4 | 4 | 6)$. Gib eine Gleichung der Geraden $h$ durch $C$ und $S$ an und begründe, dass $h$ windschief zur Geraden $BD$ ist. \hfill (Abitur 2021 GK)
\teil Ein Quader hat die Ecken $P(3 | 0 | 0)$, $R(0 | 5 | 0)$, $T(0 | 0 | 3)$ und $U(3 | 0 | 3)$. Gib eine Gleichung der Geraden $h$ durch $R$ und $U$ an und begründe, dass $h$ windschief zur Geraden $PT$ ist. \hfill (Abitur 2021 GK)
\teil Die Punkte $A(1 | 1 | 0)$, $B(3 | 3 | 4)$, $C(0 | 0 | 6)$ und $D(2 | 2 | 4)$ liegen in der Ebene $E\colon x_1 - x_2 = 0$. Die Gerade $g$ verläuft durch $A$ und $B$, die Gerade $h$ durch $C$ und $D$. Begründe, ohne den Schnittpunkt zu berechnen, dass $g$ und $h$ sich schneiden müssen. \hfill (Abitur 2018 GK)
\teil Zwei Balken liegen auf einer waagerechten Decke: der erste von $A(0 | 1 | 2)$ nach $B(4 | 3 | 2)$, der zweite von $C(1 | 5 | 2)$ nach $D(3 | -1 | 2)$ (1 LE = 1 m). Begründe, ohne den Schnittpunkt zu berechnen, dass sich die Geraden der beiden Balken schneiden. \hfill (Abitur 2018 GK)
\teil Die Punkte $A(0 | 0 | 6)$, $B(3 | 4 | 0)$, $C(1 | 5 | 4)$ und $D(2 | -1 | 2)$ liegen in der Ebene $E\colon 2x_1 + x_3 = 6$. Begründe, ohne den Schnittpunkt zu berechnen, dass sich die Gerade durch $A$ und $B$ und die Gerade durch $C$ und $D$ schneiden. \hfill (Abitur 2018 GK)
\teil Gegeben sind $g\colon \vec x = (1 | 3 | 1) + r \cdot (2 | 0 | 1)$ und die dazu echt parallele Gerade $p\colon \vec x = (1 | 5 | 1) + r \cdot (2 | 0 | 1)$. Entscheide, ob jede Gerade, die $g$ senkrecht schneidet, auch $p$ schneidet, und begründe. \hfill (Abitur 2019 GK)
\teil Gegeben sind $g\colon \vec x = (0 | 2 | 3) + r \cdot (1 | 1 | 0)$ und die dazu echt parallele Gerade $p\colon \vec x = (0 | 2 | 5) + r \cdot (1 | 1 | 0)$. Entscheide, ob jede Gerade, die $g$ senkrecht schneidet, auch $p$ schneidet, und begründe. \hfill (Abitur 2019 GK)
\end{teile}
\end{aufgabe}

% Anlauf (ohne Prüfkennung)
\begin{aufgabe}{Sachgeraden – Anlauf}
\begin{teile}
\teil Ein Seil ist zwischen zwei Masten gespannt. Wird es im Modell durch eine Gerade oder durch eine Strecke beschrieben? \kreuz{Gerade} \\ \kreuz{Strecke}
\teil Ein Seil ist an zwei Masten in $P(0 | 0 | 8)$ und $Q(6 | 8 | 5)$ befestigt (1 LE = 1 m). Gib eine Gleichung der Seilstrecke an.
\end{teile}
\end{aufgabe}

\begin{aufgabe}{Sachgeraden – Prüfungsaufgaben}
\begin{teile}
\teil Ein unteres Seil ist geradlinig von $A(0 | 0 | 6)$ nach $B(4 | 8 | 4)$ gespannt, ein oberes verläuft entlang $h\colon \vec x = (-2 | 2 | 10) + t \cdot (3 | 0 | 0)$ (1 LE = 1 m). Ein Punkt des oberen Seils liegt senkrecht über einem Punkt des unteren. Berechne den Abstand dieser beiden Punkte. \hfill (Abitur 2022 GK) \\ \leerfeld[m]
\teil Eine Stromleitung verläuft entlang $g\colon \vec x = (2 | 0 | 3) + s \cdot (2 | 6 | 1)$, ein Tragseil darüber entlang $h\colon \vec x = (0 | 3 | 8) + t \cdot (2 | 0 | -1)$ (1 LE = 1 m). Ein Punkt von $h$ liegt senkrecht über einem Punkt von $g$. Berechne den Abstand dieser beiden Punkte. \hfill (Abitur 2022 GK) \\ \leerfeld[m]
\teil Eine Wasserleitung verläuft waagerecht entlang $g\colon \vec x = (1 | 1 | 2) + s \cdot (2 | 2 | 0)$, ein Stahlseil entlang $h\colon \vec x = (0 | 6 | 9) + t \cdot (1 | -1 | -1)$ (1 LE = 1 m). Ein Punkt des Seils liegt senkrecht über einem Punkt der Leitung. Berechne den Abstand dieser beiden Punkte. \hfill (Abitur 2022 GK) \\ \leerfeld[m]
\teil Ein Seil ist geradlinig zwischen zwei Masten gespannt; die Fußpunkte sind $F_1(0 | 0 | 0)$ und $F_2(6 | 8 | 0)$, befestigt ist es in 9 m und in $7{,}5$ m Höhe (1 LE = 1 m). Berechne die Neigung des Seils in Prozent. \hfill (Abitur 2022 GK) \\ \leerfeld \%
\teil Ein Seil ist geradlinig zwischen zwei Masten gespannt; die Fußpunkte sind $F_1(0 | 0 | 0)$ und $F_2(12 | 5 | 0)$, befestigt ist es in 10 m und in $7{,}4$ m Höhe (1 LE = 1 m). Berechne die Neigung des Seils in Prozent. \hfill (Abitur 2022 GK) \\ \leerfeld \%
\teil Ein Seil ist geradlinig zwischen zwei Masten gespannt; die Fußpunkte sind $F_1(3 | 1 | 0)$ und $F_2(11 | 16 | 0)$, befestigt ist es in 8 m und in $6{,}3$ m Höhe (1 LE = 1 m). Berechne die Neigung des Seils in Prozent. \hfill (Abitur 2022 GK) \\ \leerfeld \%
\teil Ein Sonnensegel wird mit einem geradlinig gespannten Seil befestigt. Das Seil beginnt am ersten Mast im Punkt $(0 | 0 | 3)$ und endet am zweiten Mast, der senkrecht durch $M(6 | 8 | 0)$ steht. Unterwegs liegt es auf der Kante $UV$ eines waagerechten Dachs in 4 m Höhe auf, mit $U(4 | 3 | 4)$ und $V(1 | 6 | 4)$; das Verhältnis, in dem der Auflagepunkt die Kante teilt, ist bekannt (1 LE = 1 m). Beschreibe ein Verfahren, mit dem sich der Abstand des Befestigungspunkts am zweiten Mast vom Dach berechnen lässt. \hfill (Abitur 2018 GK)
\teil Ein Seil beginnt an einer Plattform in $S(1 | 0 | 2)$, liegt auf der Oberkante $KL$ einer Mauer in 3 m Höhe auf, mit $K(8 | 2 | 3)$ und $L(4 | 6 | 3)$, und ist an einem senkrechten Pfosten durch $P(15 | 10 | 0)$ befestigt. Das Verhältnis, in dem der Auflagepunkt die Kante $KL$ teilt, ist bekannt (1 LE = 1 m). Beschreibe ein Verfahren, mit dem sich berechnen lässt, wie hoch über der Mauerkrone das Seil am Pfosten befestigt ist. \hfill (Abitur 2018 GK)
\teil Eine Kamera gleitet auf einer Schiene von $A(1 | 0 | 6)$ nach $B(5 | 2 | 4)$ und filmt stets in Richtung $(2 | 1 | -3)$ (1 LE = 1 m). Um zu prüfen, ob sie den Punkt $F(8 | 3 | 0{,}5)$ filmen kann, wird der Ansatz $\vec{OK_t} + r \cdot (2 | 1 | -3) = (8 | 3 | 0{,}5)$ mit $\vec{OK_t} = (1 | 0 | 6) + t \cdot (4 | 2 | -2)$, $0 \le t \le 1$, aufgestellt. Erläutere die geometrischen Sachverhalte des Ansatzes und deute ihn im Sachzusammenhang. \hfill (Abitur 2025 GK)
\teil Ein Scheinwerfer fährt auf einer Schiene von $P(0 | 2 | 7)$ nach $Q(6 | 2 | 5)$ und strahlt stets in Richtung $(1 | 1 | -2)$ (1 LE = 1 m). Um zu prüfen, ob der Punkt $Z(7 | 6 | 0)$ auf der Bühne beleuchtet werden kann, wird der Ansatz $(0 | 2 | 7) + t \cdot (6 | 0 | -2) + r \cdot (1 | 1 | -2) = (7 | 6 | 0)$ mit $t \in [0; 1]$ aufgestellt. Erläutere die Bestandteile des Ansatzes und deute ihn im Sachzusammenhang. \hfill (Abitur 2025 GK)
\end{teile}
\end{aufgabe}

\begin{aufgabe}{Horizontalabstand aus vorgegebener Neigung in Prozent und Höhendifferenz berechnen – Prüfungsaufgaben}
\begin{teile}
\teil Ein Steg beginnt am rechten Ufer in $0{,}9$ m Höhe über dem Wasserspiegel und endet links auf der Uferzone in $0{,}3$ m Höhe; er soll eine Steigung von 4\,\% haben. Im Querschnitt liegt der Anfang bei $x = 0$, die linke Uferzone beginnt bei $x = -10$ (Angaben in m). Wie weit liegt das linke Ende des Stegs vom Rand der Uferzone entfernt? \hfill (Abitur 2017 GK) \\ \leerfeld[m]
\teil Eine Rampe mit 6\,\% Steigung soll vor einer Haustür 45 cm Höhenunterschied überwinden. Zwischen Tür und Gehweg liegen 4 m Vorgarten. Ragt die Rampe in den Gehweg hinein – und wenn ja, wie weit? \hfill (Abitur 2017 GK) \\ \leerfeld[m]
\teil Eine Rampe beginnt in $B(0 | 0 | 12)$, führt in Richtung der $x_1$-Achse mit 5\,\% Gefälle geradlinig abwärts und endet auf einer Straße in 3 m Höhe (1 LE = 1 m). Gib den Endpunkt der Rampe an. \hfill (Abitur 2017 GK)
\end{teile}
\end{aufgabe}
